\chapter{User Manual}\label{ch:ch12}

This chapter walks through APEX from the operator's chair rather than the
engineer's. Where earlier chapters explained \emph{how} the four-layer Go
backend, the Judge0 grading pipeline, and the React single-page application are
built, this one explains \emph{how to use} the finished system to run a real
classroom exam. It is organised around the two — and only two — roles the
platform recognises: \textbf{teacher} and \textbf{student}. There is no
administrator role and no separate configuration console; every operational
action is performed through the same web interface documented here.

The screen captures were taken from a local instance brought up with
\texttt{./start.sh}, which starts PostgreSQL in Docker, the Go binary on
\texttt{:8080}, and the Vite dev server on \texttt{:5173}
(\cref{fig:deploy} shows the corresponding container topology). The seed data
set provides one teacher account, one student account, one demonstration class,
and one exam of each supported question type (coding, MCQ, and written) so that
every screen below has representative content. Browser chrome has been cropped
from the captures; the routes named in the text are the real
\texttt{react-router-dom} paths served by the SPA.

\section{Getting Started: Accounts and Sign-in}

\subsection{The landing page}

An unauthenticated visitor arrives at the marketing landing page served at
\texttt{/} (\cref{fig:sc-landing}). It summarises the product — mixed-format
exams, automatic code grading, and per-test feedback — and offers a single
call to action, \emph{Get started}, which routes to the authentication screen
at \texttt{/auth}. No functionality beyond sign-up and sign-in is reachable
without an account; every \texttt{/dashboard/*} route is guarded on the client
by an auth check and on the server by the JWT authentication gate
(\cref{fig:authz}).

\begin{figure}[H]
  \centering
  
\includegraphics[width=0.85\linewidth]{F9.1_landing.png}
  \caption{The public landing page at \texttt{/}.}
  \label{fig:sc-landing}
\end{figure}

\subsection{Signing up and signing in}

The authentication screen (\cref{fig:sc-auth}) exposes three paths. A returning
user signs in with email and password; the credentials are posted to the
backend, which verifies the bcrypt hash (cost~10) and, on success, returns a
signed HS256 JSON Web Token~\cite{rfc7519,owasp-jwt} carrying the user's
identifier, email, role, and name with a 24-hour expiry. A new user chooses
\emph{Sign up}, which additionally prompts for a role — \emph{teacher} or
\emph{student} — because that single choice determines every subsequent screen
the account will ever see. A third button, \emph{Continue with Google}, delegates
to Google Identity Services~\cite{gids,oidc-core}: the browser obtains a Google
credential and the backend verifies it against the configured client ID. If the
Google account is new to APEX and no role was supplied, the backend replies that
a role is still needed and the client re-prompts before completing registration.

\begin{figure}[H]
  \centering
  \includegraphics[width=0.85\linewidth]{F9.2_auth-signin.png}
  \caption{The sign-in and sign-up screen at \texttt{/auth}. The role picker is
  shown for new sign-ups and for first-time Google accounts.}
  \label{fig:sc-auth}
\end{figure}

Passwords must be at least six characters; the email field is validated against
a standard address pattern before the request is sent. A failed login returns a
deliberately generic ``invalid email or password'' message so that the response
does not reveal whether an address is registered. If a user forgets their
password, the \emph{Forgot password} link triggers a reset flow that always
reports success (to avoid disclosing which addresses exist) and, when an SMTP
server is configured, emails a single-use, one-hour reset link; on a
development instance with no mail server the link is instead written to the
server log.

\subsection{The two dashboards}

A successful sign-in stores the token in the browser and redirects to
\texttt{/dashboard}, whose layout is chosen from the role encoded in the token.
The \textbf{teacher dashboard} (\cref{fig:sc-dash-teacher}) foregrounds the
teacher's own classes and recent activity and provides entry points to the
class list, the exam list, the grading queue, and the question bank. The
\textbf{student dashboard} (\cref{fig:sc-dash-student}) instead foregrounds the
exams that matter to a learner right now: upcoming exams, exams currently within
their time window, and recently completed ones. The two roles never share a
screen; a student has no route to the exam builder, and a teacher has no route
to the take-exam surface except through the read-only preview described later.

\begin{figure}[H]
  \centering
  \includegraphics[width=0.85\linewidth]{F9.3_dashboard-teacher.png}
  \caption{Teacher dashboard: owned classes and recent activity.}
  \label{fig:sc-dash-teacher}
\end{figure}

\begin{figure}[H]
  \centering
  
\includegraphics[width=0.85\linewidth]{F9.4_dashboard-student.png}
  \caption{Student dashboard: upcoming, active, and completed exams.}
  \label{fig:sc-dash-student}
\end{figure}

\section{Teacher Guide}

The teacher's workflow forms a single pipeline: create a class, build an exam,
assign it to the class, publish it, and — once students have submitted — grade
any manual answers and announce results. The following subsections follow that
order.

\subsection{Creating and managing classes}

The \emph{Courses} area (\texttt{/dashboard/courses}) lists the teacher's owned
classes. \emph{New class} opens a dialog that collects a name and an optional
section; on save the class is created with an automatically generated,
eight-character \textbf{invite code} that is unique across the platform
(\cref{fig:sc-class-list}). This code is the sole enrolment mechanism: there is
no roster upload and no manual add-student form. The teacher shares the code
(for example on a slide or in a chat), and students join with it themselves.

\begin{figure}[H]
  \centering
  
\includegraphics[width=0.85\linewidth]{F9.5_class-list-teacher.png}
  \caption{The teacher's class list, each card showing its eight-character
  invite code.}
  \label{fig:sc-class-list}
\end{figure}

Opening a class shows its detail screen (\cref{fig:sc-class-detail}), which
gathers three things in one place: the enrolled roster, the exams assigned to
the class, and the announcements feed. From here the teacher can rename the
class, change its cover image, adjust the passing threshold, remove a member, or
post an announcement — which fans out an in-app notification to every enrolled
student. The \emph{Grades announced} toggle is the switch that later reveals
results to students; it is examined in \cref{sec:um-results}. A safeguard
prevents accidentally announcing while manual grading is still outstanding, so
that no student sees a total that is about to change.

\begin{figure}[H]
  \centering
  
\includegraphics[width=0.85\linewidth]{F9.6_class-detail-teacher.png}
  \caption{Class detail (teacher view): roster, assigned exams, and
  announcements.}
  \label{fig:sc-class-detail}
\end{figure}

\subsection{Building an exam}

\emph{Exams} (\texttt{/dashboard/exams}) lists the teacher's exams with status
badges — Draft, Upcoming, Active, or Completed — derived from the draft flag and
the start and end times. \emph{New exam} opens the \textbf{exam builder}
(\cref{fig:sc-builder}). The builder is a three-part surface: a header collects
exam-level metadata (title, description, duration in minutes, start and end
times, and behavioural toggles such as \emph{shuffle questions} and \emph{show
results after}); a left rail lists the questions in order; and the right pane
edits whichever question is selected. The same builder assembles all three
question types into a single sitting, which is the platform's central idea — a
student answers coding, multiple-choice, and written questions in one timed
attempt, and one grading pipeline stands behind all three.

\begin{figure}[H]
  \centering
  
\includegraphics[width=0.85\linewidth]{F9.7_exam-builder.png}
  \caption{The exam builder: metadata header, question rail, and editing pane.}
  \label{fig:sc-builder}
\end{figure}

\subsubsection{Coding questions}

A coding question (\cref{fig:sc-coding-editor}) carries a title, a Markdown
problem statement, optional starter code, and — most importantly — a list of
\textbf{test cases}. Each test case pairs an input with an expected output and
is flagged as either \emph{sample} or \emph{hidden}. Sample test cases are shown
to the student and can be run during the attempt for immediate feedback; hidden
test cases are withheld and are executed only when the exam is graded, which is
what prevents a student from reverse-engineering the full grader. Each question
also declares an execution time limit (default \texttt{2000}\,ms) and a memory
limit (default \texttt{262144}\,KB, i.e.\ 256\,MB) that are passed to the Judge0
sandbox~\cite{judge0,isolate} when the code runs. The scoring is proportional:
the earned score is the fraction of test cases that pass, so a solution that
clears three of four tests scores 75\%.

\begin{figure}[H]
  \centering
  
\includegraphics[width=0.85\linewidth]{F9.8_problem-editor-coding.png}
  \caption{Coding-question editor: statement, starter code, limits, and
  sample/hidden test cases.}
  \label{fig:sc-coding-editor}
\end{figure}

\subsubsection{Multiple-choice questions}

A multiple-choice question (\cref{fig:sc-mcq-editor}) is a list of options, each
with a \emph{Correct} checkbox. Marking more than one option correct turns the
question into a multiple-answer question, graded by exact set equality: the
student earns full marks only if their selected set matches the correct set
exactly, otherwise zero. The editor requires at least one option to be marked
correct — if a question is published with no correct answer configured, the
grader cannot fairly fail anyone on it and instead routes that submission to the
manual-review queue rather than awarding a wrong-answer zero.

\begin{figure}[H]
  \centering
  
\includegraphics[width=0.85\linewidth]{F9.9_problem-editor-mcq.png}
  \caption{Multiple-choice editor: options with per-option \emph{Correct}
  checkboxes.}
  \label{fig:sc-mcq-editor}
\end{figure}

\subsubsection{Written questions}

A written question (\cref{fig:sc-written-editor}) collects a prompt, an optional
grading rubric, and a maximum word count (default 500). Written answers are
never graded automatically — every written submission begins life in the
\emph{pending review} state and enters the manual-grading queue, where the
teacher reads it and assigns a score. Crucially, a written answer that is still
awaiting a manual grade is excluded from the student's running total, so an
in-flight essay grade never artificially depresses the score the student sees.

\begin{figure}[H]
  \centering
  
\includegraphics[width=0.85\linewidth]{F9.10_problem-editor-written.png}
  \caption{Written-question editor: prompt, rubric, and word-count limit.}
  \label{fig:sc-written-editor}
\end{figure}

\Cref{tab:um-qtypes} summarises how the three question types differ from the
teacher's point of view.

\begin{table}[H]
  \centering
  \caption{The three question types and how each is graded.}
  \label{tab:um-qtypes}
  \small
  \begin{tabularx}{\linewidth}{@{}L{2.4cm} L{3.4cm} X@{}}
    \toprule
    \textbf{Type} & \textbf{Author supplies} & \textbf{How it is scored} \\
    \midrule
    Coding &
      Statement, starter code, time/memory limits, sample and hidden test cases &
      Automatically: code runs in the Judge0 sandbox against every test case;
      score is the fraction of tests passed. \\
    \addlinespace
    Multiple choice &
      Options with one or more correct answers, optional explanation &
      Automatically: exact set-equality on selected vs.\ correct options; all or
      nothing. Unconfigured answers go to manual review. \\
    \addlinespace
    Written &
      Prompt, optional rubric, maximum word count &
      Manually: always starts in \emph{pending review}; the teacher assigns a
      score from the grading queue. \\
    \bottomrule
  \end{tabularx}
\end{table}

\subsection{The question bank}

Authoring good questions is expensive, so a teacher rarely wants to write each
one from scratch for every exam. The \textbf{question bank}
(\texttt{/dashboard/question-bank}, teacher only) lets the teacher save
questions independently of any exam, organise them into folders, and later
import them into a new exam. Bank questions and exam-attached questions share the
same underlying storage but are distinguished by a bank flag, so importing a
question copies it into the exam without disturbing the reusable original.

\subsection{Assigning, publishing, and previewing}

Building a draft exam does not expose it to anyone. Two further steps are
required. First, the exam must be \textbf{assigned} to one or more classes; the
assignment step validates that every target class actually belongs to the
teacher performing it, so a teacher cannot push an exam into a colleague's
class. Second, the exam must be \textbf{published}. The \emph{Publish} toggle is
enabled only once the exam has at least one question and a non-empty start time
(\cref{fig:sc-publish}); an exam with no questions or no schedule stays a draft.
This gate matters because student-facing queries deliberately hide draft exams —
the same filter whose accidental corruption is the ``\texttt{is\_draft}'' war
story recounted in \cref{fig:migration}.

\begin{figure}[H]
  \centering
  \includegraphics[width=0.85\linewidth]{F9.11_publish-toggle.png}
  \caption{The publish toggle, enabled once the exam has questions and a start
  time.}
  \label{fig:sc-publish}
\end{figure}

Before publishing, the teacher can \emph{Preview} the exam
(\cref{fig:sc-preview}), which renders it exactly as a student will see it —
question rail, Monaco editor, and all — but read-only and without creating an
attempt. Previewing is the recommended last check that starter code compiles,
sample tests behave, and the wording is unambiguous.

\begin{figure}[H]
  \centering
  
\includegraphics[width=0.85\linewidth]{F9.14_exam-preview.png}
  \caption{Exam preview: the student-facing view, rendered read-only for the
  teacher.}
  \label{fig:sc-preview}
\end{figure}

\section{Student Guide}

The student's workflow is the mirror image of the teacher's: join a class, wait
for an exam to open, take it within its window, and later review the result.

\subsection{Joining a class}

A student opens \emph{Courses} and chooses \emph{Join with code}
(\cref{fig:sc-join}). The eight-character invite code is the only input.
Enrolment is idempotent — a student can be a member of a given class at most
once — so re-entering the same code is harmless. Once joined, the class appears
on the dashboard, and any exam already assigned to that class surfaces under
\emph{Upcoming}.

\begin{figure}[H]
  \centering
  \includegraphics[width=0.85\linewidth]{F9.15_join-class.png}
  \caption{Joining a class with its invite code.}
  \label{fig:sc-join}
\end{figure}

\subsection{Taking a timed exam}

When an exam is within its window, \emph{Start exam} opens the take screen at
\texttt{/dashboard/exam/:id} (\cref{fig:sc-take}). The layout is consistent
across question types: a left rail numbers the questions and shows which are
answered, the remaining time counts down in the top corner, and a Submit button
sits at the end. Coding questions render in the Monaco editor~\cite{monaco} —
the same engine that powers Visual Studio Code — with syntax highlighting and a
language selector covering the six supported languages (Python, JavaScript,
TypeScript, Java, C, and C++). Multiple-choice questions render as a checkbox
list, and written questions render as a textarea with a live word counter that
enforces the author's limit.

\begin{figure}[H]
  \centering
  \includegraphics[width=0.85\linewidth]{F9.16_take-coding.png}
  \caption{Taking a coding question in the Monaco editor, with the question rail
  and countdown timer.}
  \label{fig:sc-take}
\end{figure}

\paragraph{Running sample tests.} On a coding question the student can run their
code against the \emph{sample} test cases at any time before submitting. This
sends the current code and chosen language to the backend, which executes it in
the Judge0 sandbox and returns per-test results — expected versus actual output
— rendered inline. Only sample tests run on this path; the hidden tests that
ultimately decide the grade are never revealed here, and a sample run does not
count as a submission.

\paragraph{Autosave.} The client periodically saves a draft of every answer to
the server while the attempt is open. Because the draft lives on the server and
not merely in the browser, a crash, an accidental tab close, or a switch to
another device does not lose work: reopening the exam restores the saved
answers. Output comparison during grading is whitespace-tolerant — trailing
spaces and Windows versus Unix line endings are normalised on both sides — so a
correct answer is not marked wrong over an invisible newline difference.

\paragraph{Submitting.} \emph{Submit} raises a confirmation dialog
(\cref{fig:sc-submit-dialog}) — submission is final and ends the attempt, so the
extra click is deliberate. On confirmation the client sends every answer to the
server, which records the submissions atomically and immediately begins grading
them. The student then sees the post-submission screen
(\cref{fig:sc-postsubmit}), which acknowledges receipt and explains that
automatically graded questions score right away while any written answers await
the teacher.

\begin{figure}[H]
  \centering
  \includegraphics[width=0.85\linewidth]{F9.19_submit-dialog.png}
  \caption{The submit confirmation dialog.}
  \label{fig:sc-submit-dialog}
\end{figure}

\begin{figure}[H]
  \centering
  \includegraphics[width=0.85\linewidth]{F9.20_post-submit.png}
  \caption{The post-submission acknowledgement screen.}
  \label{fig:sc-postsubmit}
\end{figure}

\section{Grading and Results}\label{sec:um-results}

\subsection{The manual-grading queue}

Coding and multiple-choice answers are scored automatically within seconds of
submission, but written answers wait for a human. The teacher's \emph{Grading}
area (\texttt{/dashboard/grading}, \cref{fig:sc-grading-queue}) lists every
submission still in \emph{pending review}. Opening one shows the student's
answer alongside the rubric and inputs for a numeric score and free-text
feedback (\cref{fig:sc-grading-detail}). Saving a grade both records the score
and re-aggregates the student's overall exam total; the sequence of steps from
submission through automatic grading to this manual step is diagrammed in
\cref{fig:grading-seq}. Only the teacher who owns the exam may grade its
submissions — the ownership check is the third authorization gate in
\cref{fig:authz}.

\begin{figure}[H]
  \centering
  
\includegraphics[width=0.85\linewidth]{F9.21_grading-queue.png}
  \caption{The pending-grading queue of written answers awaiting a score.}
  \label{fig:sc-grading-queue}
\end{figure}

\begin{figure}[H]
  \centering
  
\includegraphics[width=0.85\linewidth]{F9.22_grading-detail.png}
  \caption{Grading a written answer against its rubric, with score and feedback
  inputs.}
  \label{fig:sc-grading-detail}
\end{figure}

\subsection{Reviewing results as a teacher}

The per-exam results screen (\cref{fig:sc-results-student}) presents one row per
student — total score, percentage, and time taken — giving the teacher an
at-a-glance view of how a class performed. Selecting a row drills into that
student's per-question breakdown (\cref{fig:sc-results-detail}), which shows how
each answer scored, including per-test outcomes for coding questions. This is
the teacher's audit trail: it is where a disputed grade can be traced back to the
exact test case that failed.

\begin{figure}[H]
  \centering
  \includegraphics[width=0.85\linewidth]{F9.23_results-by-student.png}
  \caption{Per-student results for one exam: score, percentage, and time taken.}
  \label{fig:sc-results-student}
\end{figure}

\begin{figure}[H]
  \centering
  \includegraphics[width=0.85\linewidth]{F9.24_results-detail.png}
  \caption{Per-question breakdown for a single student's attempt.}
  \label{fig:sc-results-detail}
\end{figure}

\subsection{Announcing and reviewing results as a student}

Results are not visible to students the instant grading finishes. Two conditions
must hold: the exam's \emph{show results after} behaviour must be on, and the
teacher must flip \emph{Grades announced} for the class. This two-step gate lets
a teacher finish manual grading, sanity-check the distribution, and then release
every student's result together. Once released, the student's \emph{Results} area
(\cref{fig:sc-review}) lists their attempts; opening one reveals the per-question
result, the teacher's written feedback, and — for coding questions — the
outcomes of the sample tests only, keeping the hidden tests confidential for
reuse in future exams.

\begin{figure}[H]
  \centering
  \includegraphics[width=0.85\linewidth]{F9.25_student-review.png}
  \caption{The student's review screen: per-question results and teacher
  feedback.}
  \label{fig:sc-review}
\end{figure}

\section{Notifications, Profile, and Settings}

Beyond the exam pipeline, three supporting surfaces are shared by both roles.

\paragraph{Notifications.} A bell icon in the top bar opens a popover of recent
notifications, and a dedicated notifications page lists the full history. Each
notification carries a deep link to the thing it concerns — a new assignment, a
posted announcement, an imminent exam, or the ``Exam Graded'' event that fires
once a student's whole attempt has finished grading. Exam reminders are also
delivered by email: the in-process scheduler sweeps for published exams and
sends both an in-app and an email reminder roughly an hour before an exam
starts, and a second email as it begins, honouring each user's notification
preferences.

\paragraph{Profile.} The profile page lets a user edit their display name, a
short biography, and an avatar image.

\paragraph{Settings.} The settings page governs how the platform contacts the
user. Separate toggles control in-app versus email notifications for each
category — announcements, results, and the exam-reminder emails that the
scheduler respects. Teachers additionally see per-account defaults, such as the
default passing threshold applied to new exams and whether new exams start as
drafts. This page is also where a user can permanently delete their account,
which cascades through their owned or enrolled data in a single transaction.

\section{Summary}

APEX presents two coherent, role-shaped paths through one shared data model. The
teacher path runs \emph{create class $\rightarrow$ build mixed-format exam
$\rightarrow$ assign and publish $\rightarrow$ grade $\rightarrow$ announce}; the
student path runs \emph{join by code $\rightarrow$ take timed exam in Monaco
$\rightarrow$ submit $\rightarrow$ review}. The screenshots in this chapter are
the load-bearing user-facing artefacts of the project, and together they show
that the architecture described in the preceding chapters resolves, at the
surface, into a small number of focused screens that a working teacher and a
working student can operate without training.
